% GENERATED FILE. Edit canonical lessons, then run npm run generate:books.
% canonical-source-hash: fnv1a64:7959698a4626594a
% canonical-lessons: ZH-C384-zhongshu, ZH-C384-chengwei, ZH-W384-cheng-become, ZH-C384-chengben, ZH-C384-ziyuan, ZH-C384-zengjia

\chapter{Becoming Green}
\label{ch:zh-becoming-green}
\hlchaptermodality{384}

\begin{chapteropening}
\textbf{By the end of this chapter:} \emph{I can talk about planting trees, becoming something, what things cost to make, resources and growth.}

Complete the last lesson of chapter 384: I can talk about planting trees, becoming something, what things cost to make, resources and growth.
\end{chapteropening}

\section[zhòngshù]{zhòngshù — to plant trees}

\label{lesson:ZH-C384-zhongshu}



\begin{quote}
Before the new one: say the Mandarin for to cut down on, then the Mandarin for disposable (used once).
\end{quote}

\subsection*{You'll want to know: zhòngshù}
\emph{zhòngshù} — \textquotedblleft{}to plant trees\textquotedblright{}.

\begin{quote}
\emph{měi nián chūntiān wǒmen dōu qù shān lǐ zhòngshù} — Every spring we go to the hills to plant trees.
\end{quote}

Its characters come later; for now it is a word for the ear.

\begin{cousinweb}
\emph{zhòngshù} is built from \emph{zhòng} 'to plant' and \emph{shù} 'a tree', both inherited from Middle Chinese, \emph{tsyowngH} and \emph{dzyuH} (Guangyun, in Baxter's transcription), which glosses the second 'trees as a whole' and 'to set up'.
\end{cousinweb}

\begin{grammarlens}[title={What you've built}]
One more word for giving reasons.
\end{grammarlens}

\subsection*{Your turn}
\begin{itemize}
\raggedright
  \item \emph{Say it:} \emph{zhòngshù}
  \item \emph{Say it:} \emph{zhòngshù}, once more
  \item \emph{Say it:} \emph{měi nián chūntiān wǒmen dōu qù shān lǐ zhòngshù} — Every spring we go to the hills to plant trees.
  \item \emph{Recall:} say the Mandarin for to cut down on, then the Mandarin for disposable (used once), then say \emph{zhòngshù} again
\end{itemize}

\begin{tcolorbox}[breakable,colback=teal!4,colframe=teal!35!black,title={Before you move on}]
What does zhòngshù mean? (\textquotedblleft{}To plant trees\textquotedblright{}.) Say it once more.
\end{tcolorbox}
\section[\texorpdfstring{\zh{成为} (chéngwéi)}{成为 (chéngwéi)}]{\zh{成为} (chéngwéi) — to become}

\label{lesson:ZH-C384-chengwei}



\begin{quote}
Before the new one: say the Mandarin for disposable (used once), then the Mandarin for to plant trees.
\end{quote}

\subsection*{You'll want to know: \zh{成为}}
\textbf{\zh{成为}} — \emph{chéngwéi} — \textquotedblleft{}to become\textquotedblright{}.

\begin{quote}
\textbf{\zh{他想成为一个医生}} \emph{tā xiǎng chéngwéi yí ge yīshēng} — He wants to become a doctor.
\end{quote}

The new character, in pieces: \zh{成} — six strokes: a short bar, a long stroke falling to the left, a small turning hook, the long slanting hook of \zh{我}, a short fall and a dot. The next lesson puts it in your hand.

\begin{cousinweb}
\textbf{\zh{成为}} (\emph{chéngwéi}) is built from \zh{成} 'to complete' and \zh{为} 'to make', both inherited from Middle Chinese, \emph{dzyeng} and \emph{hjwe} (Guangyun, in Baxter's transcription), which glosses the first 'to finish' and the second 'to do'.
\end{cousinweb}

\begin{grammarlens}[title={What you've built}]
One more word for giving reasons.
\end{grammarlens}

\subsection*{Your turn}
\begin{itemize}
\raggedright
  \item \emph{Say it:} \emph{chéngwéi}
  \item \emph{Say it:} \emph{chéngwéi}, once more
  \item \emph{Say it:} \emph{tā xiǎng chéngwéi yí ge yīshēng} — He wants to become a doctor.
  \item \emph{Recall:} say the Mandarin for disposable (used once), then the Mandarin for to plant trees, then say \emph{chéngwéi} again
\end{itemize}

\begin{tcolorbox}[breakable,colback=teal!4,colframe=teal!35!black,title={Before you move on}]
What does \zh{成为} mean? (\textquotedblleft{}To become\textquotedblright{}.) Say it once more.
\end{tcolorbox}
\section[\texorpdfstring{\zh{成} (chéng)}{成 (chéng)}]{\zh{成} — a bar, a fall, a small turn, then the slanting hook of \zh{我}}

\label{lesson:ZH-W384-cheng-become}



\begin{quote}
The word first, then the character inside it.

\begin{itemize}
\raggedright
  \item \emph{Recall:} say \emph{chéngwéi} — \textbf{R1}, one lesson back, before it decays
\end{itemize}

\begin{quote}
\zh{我} \zh{为}
\end{quote}

\emph{Read it:} each one aloud, with nothing said first — \emph{wǒ}, \emph{wéi}
\end{quote}

\subsection*{Script — a bar, a fall, a small turn, then the slanting hook of \zh{我}}
\hlblockfigure{\detokenize{figures/ZH-W384-cheng-become-filmstrip.pdf}}{How \zh{成} is written, stroke by stroke}

\begin{quote}
\zh{成}
\end{quote}

\textbf{One shape, six strokes in a fixed order.} The order is the thing to hold.

\textbf{A bar,} left to right, from the left edge to past the middle.

\textbf{A long stroke falling to the left} from the left end of the bar down to the foot.

\textbf{A small turning hook} under the bar: a short bar that turns down and hooks to the left.

\textbf{The long slanting hook,} as in \zh{我}: from the top, a long stroke down to the lower right that flicks up at the end. Then \textbf{a short stroke falling to the left} across it, and \textbf{a dot} at the upper right.

Six strokes in all.

\subsection*{Your turn}
\begin{itemize}
\raggedright
  \item \emph{Write it:} \zh{成}, all six, with the model in front of you
  \item \emph{Cover:} write it again from memory, and count the pieces, not the strokes
\end{itemize}

\begin{tcolorbox}[breakable,colback=teal!4,colframe=teal!35!black,title={Before you move on}]
How many strokes in \zh{成}? (\textbf{Six}.) Which stroke flicks up at the end? (\textbf{The long slanting hook, the fourth}.) What comes last? (\textbf{The dot at the upper right}.)

Source: \href{https://raw.githubusercontent.com/chanind/hanzi-writer-data/68d10a4b21150cae5e1ebbd223eed289cf32d90c/data/\%E6\%88\%90.json}{Hanzi Writer Data \zh{成}, pinned stroke paths and medians}.
\end{tcolorbox}
\section[\texorpdfstring{\zh{成本} (chéngběn)}{成本 (chéngběn)}]{\zh{成本} (chéngběn) — production cost; cost}

\label{lesson:ZH-C384-chengben}



\begin{quote}
Before the new one: say the Mandarin for to plant trees, then the Mandarin for to become.
\end{quote}

\subsection*{You'll want to know: \zh{成本}}
\textbf{\zh{成本}} — \emph{chéngběn} — \textquotedblleft{}production cost; cost\textquotedblright{}.

\begin{quote}
\textbf{\zh{请保安的成本很高}} \emph{qǐng bǎo'ān de chéngběn hěn gāo} — It costs a lot to hire security guards.
\end{quote}

\begin{cousinweb}
\textbf{\zh{成本}} (\emph{chéngběn}) is built from \zh{成} 'to complete' and \zh{本} 'the root', both inherited from Middle Chinese, \emph{dzyeng} and \emph{pwonX} (Guangyun, in Baxter's transcription), which glosses the second 'the foot of a tree': what a finished thing rests on. The sense 'cost' is not given there.
\end{cousinweb}

\begin{grammarlens}[title={What you've built}]
One more word for giving reasons.
\end{grammarlens}

\subsection*{Your turn}
\begin{itemize}
\raggedright
  \item \emph{Say it:} \emph{chéngběn}
  \item \emph{Say it:} \emph{chéngběn}, once more
  \item \emph{Say it:} \emph{qǐng bǎo'ān de chéngběn hěn gāo} — It costs a lot to hire security guards.
  \item \emph{Recall:} say the Mandarin for to plant trees, then the Mandarin for to become, then say \emph{chéngběn} again
\end{itemize}

\begin{tcolorbox}[breakable,colback=teal!4,colframe=teal!35!black,title={Before you move on}]
What does \zh{成本} mean? (\textquotedblleft{}Production cost; cost\textquotedblright{}.) Say it once more.
\end{tcolorbox}
\section[zīyuán]{zīyuán — natural resources}

\label{lesson:ZH-C384-ziyuan}



\begin{quote}
Before the new one: say the Mandarin for to become, then the Mandarin for production cost.
\end{quote}

\subsection*{You'll want to know: zīyuán}
\emph{zīyuán} — \textquotedblleft{}natural resources\textquotedblright{}.

\begin{quote}
\emph{shuǐ shì hěn zhòngyào de zīyuán suǒyǐ wǒmen bù yīnggāi làngfèi} — Water is a very important resource, so we shouldn't waste it.
\end{quote}

Its characters come later; for now it is a word for the ear.

\begin{cousinweb}
\emph{zīyuán} is built from \emph{zī} 'goods' and \emph{yuán} 'a source', both inherited from Middle Chinese, \emph{tsij} and \emph{ngjwon} (Guangyun, in Baxter's transcription), which glosses the first 'wealth' and 'to help' and the second 'where water rises'.
\end{cousinweb}

\begin{grammarlens}[title={What you've built}]
One more word for giving reasons.
\end{grammarlens}

\subsection*{Your turn}
\begin{itemize}
\raggedright
  \item \emph{Say it:} \emph{zīyuán}
  \item \emph{Say it:} \emph{zīyuán}, once more
  \item \emph{Say it:} \emph{shuǐ shì hěn zhòngyào de zīyuán suǒyǐ wǒmen bù yīnggāi làngfèi} — Water is a very important resource, so we shouldn't waste it.
  \item \emph{Recall:} say the Mandarin for to become, then the Mandarin for production cost, then say \emph{zīyuán} again
\end{itemize}

\begin{tcolorbox}[breakable,colback=teal!4,colframe=teal!35!black,title={Before you move on}]
What does zīyuán mean? (\textquotedblleft{}Natural resources\textquotedblright{}.) Say it once more.
\end{tcolorbox}
\section[zēngjiā]{zēngjiā — to increase}

\label{lesson:ZH-C384-zengjia}



\begin{quote}
Before the new one: say the Mandarin for production cost, then the Mandarin for natural resources.
\end{quote}

\subsection*{You'll want to know: zēngjiā}
\emph{zēngjiā} — \textquotedblleft{}to increase\textquotedblright{}.

\begin{quote}
\emph{yīnwèi rénkǒu zēngjiā le suǒyǐ fángzi hěn guì} — The population has grown, so houses are expensive.
\end{quote}

Its characters come later; for now it is a word for the ear.

\begin{cousinweb}
\emph{zēngjiā} is built from \emph{zēng} and \emph{jiā}, both inherited from Middle Chinese, \emph{tsong} and \emph{kae} (Guangyun, in Baxter's transcription), which glosses each with the other: 'to add'.
\end{cousinweb}

\begin{grammarlens}[title={What you've built}]
One more word for giving reasons.
\end{grammarlens}

\subsection*{Your turn}
\begin{itemize}
\raggedright
  \item \emph{Say it:} \emph{zēngjiā}
  \item \emph{Say it:} \emph{zēngjiā}, once more
  \item \emph{Say it:} \emph{yīnwèi rénkǒu zēngjiā le suǒyǐ fángzi hěn guì} — The population has grown, so houses are expensive.
  \item \emph{Recall:} say the Mandarin for production cost, then the Mandarin for natural resources, then say \emph{zēngjiā} again
\end{itemize}

\begin{tcolorbox}[breakable,colback=teal!4,colframe=teal!35!black,title={Before you move on}]
What does zēngjiā mean? (\textquotedblleft{}To increase\textquotedblright{}.) Say it once more.
\end{tcolorbox}
